% ------------------------------------------------------------- Discussion
\paragraph{What a defender should take away.}
Four design rules follow from the measurements, and each is a mitigation of
the attack the paper describes.
\begin{enumerate}
\item \emph{Re-anchor, and size the interval rather than the threshold.} With a
  trusted fix every 5\,s and its gate armed, frequent evasions contributed at
  most \nGatedCap\,m and the pooled supremum was \nAnchorMaxFive\,m, for every
  detector and disguise. Without the gate, the contribution reached
  \nNgMax\,m. Lengthening the interval from 5 to 20\,s multiplied the budget by
  \nAnchorMultTwenty{} (\cref{fig:anchor}), while moving the operating point
  from 1\% to 10\% FPR left the pooled budget unchanged (\cref{sec:eval}). For a corridor of
  margin $m$, choose $T_a$ such that
  $\dadv(T_a)\le m-e^{LT_c}\dnet-s^\star$, where $s^\star$ is the smallest
  budget with $\pbar(s^\star)$ below the tolerated risk.
\item \emph{Close the last interval.} An unchecked segment before mission end
  lets the attacker skip the return leg and multiplied the budget by
  \nOpenRatioTwenty{} at $T_a=20$\,s. Landing and terminal approach are where
  a trusted fix matters most.
\item \emph{Model the inertial bias, not just cross-check the inertial sensor.}
  Without the anchor gate, the only detector that capped every disguise was the
  one whose filter models the accelerometer bias (D2). D1 and D3 admitted tens
  of metres to the GNSS and weather mimics, and even the inertial invariant D4
  admitted \nNgCapDFour\,m to the fault mimic, whose bias-shaped curvature it
  tolerates by design.
\item \emph{Certify the undetected outcome, not the alarm.} \Cref{thm:floor}
  bounds $\Pr[\text{fail}\wedge U]$ for every attacker whose undetected
  contribution stays within $\dadv$, without relying on the detector catching
  any particular disguise. A floor computed from benign marginals alone and
  gated on silence has no such property (\cref{prop:gated}).
\end{enumerate}

\paragraph{On the separability counter-claim.}
Li et al.~\cite{li2026unified} and Xue et al.~\cite{xue2025dual} argue that
faults and attacks act through different subspaces and can be separated by
residual design. Our gate-free results are compatible with that view for a
\emph{given} residual: D2 and D4 do separate several disguises. They differ in
the quantifier. Separability holds for a fixed attack against a fixed detector.
Equivalence concerns the attacker's best response, and against every detector
we built, some family of small-amplitude offsets evaded in at least
\nCrossEvMin{} of missions. Whether a residual exists that tolerates the benign
classes and still rejects all of their imitations is open. Our evidence says
only that none of ours does.

\paragraph{Limitations.}
\begin{itemize}
\item \emph{Reduced-order model.} The pilot simulates one lateral channel with
  a linear KF. PX4's EKF2 has more states, different gating and nonlinear
  coupling. The disguise matrix and soundness must be re-established in SITL
  (runbook R3) before they are claimed for a deployed autopilot.
\item \emph{Empirical adversarial term.} $\dadv$ is a pooled supremum over
  \nPooledAttackers{} CEM attackers, not a proof. The held-out attacker exceeded
  it in \nSupExceed{} of \nSupChecks{} cells by up to \nHeldoutAdvRatio$\times$.
  Stronger optimisers (R1) and an analytic stealthy-reachable-set
  bound~\cite{mo2016,murguia2020} (R8) are needed.
\item \emph{Conservative bound, limited power.} The observed frequency of
  undetected failure was zero wherever the bound was informative, so
  soundness was tested mainly through negative controls. Closing the temporal
  gap would make the bound tight enough to be stressed directly.
\item \emph{Simulated marginals.} The benign marginals come from the
  simulator; the runbook replaces them with the Whelan campaign and the RflyMAD
  and ALFA flights (R5).
\item \emph{Network term not simulated.} $\dnet$ enters as the parent's
  certified tube; R4 injects measured relay staleness into SITL.
\item \emph{Covert attackers only.} The bound covers undetected failure. An
  attacker indifferent to detection is bounded instead by detection latency and
  the failsafe, which we do not certify here.
\item \emph{No IMU injection.} An attacker who can also inject into the IMU
  (acoustically~\cite{son2015,trippel2017} or through bus
  reconfiguration~\cite{erba2025}) defeats the inertial cross-checks, and
  condition C2$'$ then carries the whole burden.
\end{itemize}
